%% Article 2 (workshop): the offline half only.
%%
%% Target venue: a workshop on LLM-assisted constraint acquisition. CONFIRM
%% IT IS NON-ARCHIVAL, or permits a later journal version, before submitting:
%% an archival workshop paper would complicate the main article.
%%
%% Build: latexmk -xelatex workshop.tex
%%
%% DRAFT STATUS. Results sections carry their heading and nothing else.

\documentclass[sigconf,screen,review,anonymous]{acmart}

\usepackage{booktabs}

\newcommand{\result}[1]{\textbf{[\textsc{result}: #1]}}
\newcommand{\awaitingdata}[1]{%
  \par\noindent\textit{[Awaiting data: #1]}\par}

\settopmatter{printacmref=false}
\setcopyright{none}

\begin{document}

\title{Formalising Staff-Authored Timetabling Constraints}
\subtitle{A Labelled Corpus and a Cross-Model Accuracy Study}

\author{Anonymous Author(s)}
\affiliation{%
  \institution{Anonymous Institution}
  \country{}
}

\begin{abstract}
Benchmarks for turning natural language into formal constraints have grown
quickly, and scheduling is now among the domains they cover. Their natural
language, however, is rendered from templates: in the leading scheduling
benchmark no domain user wrote any of it, and its own limitations section
names human-authored phrasing as the route to better ecological validity.
Whether a model that handles templated prose handles an email from a lecturer
is an open question.

We release a corpus of \result{K} timetabling constraint sentences written by
university staff, labelled independently by two annotators against a closed
vocabulary of ten constraint types, with Cohen's $\kappa$ reported on both the
type and the strength assignment. We measure translation accuracy on it across
\result{M} models from \result{P} providers under one instrument, decomposing
errors into type, strength and entity categories, and separating entity
failures the system could not resolve from those it resolved to the wrong
entity. We also report the distribution of requests the vocabulary cannot
express, which is a property of what staff ask for and not a defect rate.
\end{abstract}

\begin{CCSXML}
<ccs2012>
<concept>
<concept_id>10010147.10010178.10010187</concept_id>
<concept_desc>Computing methodologies~Knowledge representation and reasoning</concept_desc>
<concept_significance>500</concept_significance>
</concept>
<concept>
<concept_id>10010147.10010178.10010179.10010182</concept_id>
<concept_desc>Computing methodologies~Natural language generation</concept_desc>
<concept_significance>300</concept_significance>
</concept>
</ccs2012>
\end{CCSXML}

\ccsdesc[500]{Computing methodologies~Knowledge representation and reasoning}
\ccsdesc[300]{Computing methodologies~Natural language generation}

\keywords{constraint acquisition, timetabling, large language models,
labelled corpus, inter-annotator agreement}

\maketitle

%% =====================================================================
\section{Introduction}
%% =====================================================================

Constraint modelling is the bottleneck for wider adoption of constraint
programming~\cite{michailidis2025cpbench}, and language models are now applied
to it directly~\cite{shi2025constraintllm,ramamonjison2022nl4opt}. For
scheduling, SCHEDBench puts six problem families including curriculum
timetabling to thirteen models and reports that none is invariant to how the
same problem is worded~\cite{sharma2026schedbench}.

Two properties of that evaluation leave a question open. Its natural language
comes from a deterministic template pipeline with slot filling, so no domain
user wrote any of it, and the paper itself names broader human-authored
phrasing as what it lacks. And it scores end-to-end schedule generation,
stating that its invariance claims do not isolate effects to the modelling
stage.

This paper supplies the missing input and isolates the stage. We collected
timetabling constraint sentences written by university staff, labelled them
twice, and measured how accurately models map them onto a closed constraint
vocabulary. No timetable is generated: the model's only output is constraints.

Three contributions:

\begin{itemize}
  \item A corpus of staff-authored timetabling constraint sentences with two
        independent label sets and agreement reported. Existing scheduling
        benchmarks template their language and report no agreement statistic.
  \item A cross-model, cross-provider accuracy measurement under one
        instrument, decomposed by error category.
  \item An error taxonomy grounded in what staff actually wrote, including the
        class of references that pick out an entity without naming it.
\end{itemize}

%% =====================================================================
\section{Corpus}
%% =====================================================================

\subsection{Collection}

Constraint sentences written by university teaching and administrative staff
in the course of describing their own timetabling requirements. The brief
given to contributors asks for requirements, not for sentences in any
particular form, so the corpus carries the hedging, the implicit references
and the compound sentences that real requests carry.

\subsection{The Target Vocabulary}

Ten constraint types. Eight are solvable: lecturer unavailability, lecturer
time preference, a course pinned to a named room, a course requiring a type of
room, a cap on consecutive periods, an earliest acceptable period, a pair of
courses that must not overlap, and a set of courses to be placed on one day.
Each carries a strength, hard or soft.

Two stay open. \texttt{unsupported} marks a schedulable request the vocabulary
cannot express; \texttt{needs\_clarification} marks input too ambiguous to
formalise. Neither reaches a solver, and their distribution is reported as a
finding about what staff ask for.

\subsection{Labelling}

Two annotators independently assign the constraint or constraints each
sentence expresses, using canonical entity identifiers. Agreement is reported
as Cohen's $\kappa$~\cite{cohen1960kappa} on the type assignment and on the
strength assignment. Because $\kappa$ requires one category per item, the
primary constraint of each sentence supplies the label, and the number of
sentences carrying more than one constraint is reported so the cost of that
convention is visible; set-level exact agreement uses every constraint.
Disagreements are adjudicated to a gold label.

%% =====================================================================
\section{Measurement}
%% =====================================================================

\subsection{Pipeline}

A system prompt carries the vocabulary, the department's register of entities
and sixteen worked examples. The reply is validated against the schema, with a
single retry on a validation failure. A refusal, a response truncated at the
token ceiling, and an API error are each recorded as that item's outcome and
never read as a result.

Entity references are resolved against the register after translation, which
keeps the output contract simple and makes resolution failures measurable
separately from translation failures.

\subsection{Measures}

Exact match at the sentence level; type accuracy; hard/soft accuracy
conditional on a correct type; and entity accuracy decomposed into correct,
unresolved, misresolved, missed and spurious.

Descriptive references are counted separately. We define a descriptive
reference as one that picks out an entity without naming it: ``the big lab'',
``the first years''. LlmAcq found this to be the class people actually
produce, reporting that users rely on implicit structural knowledge so that a
sentence may name no variable at all~\cite{mechqrane2026llmacq}. A description
resolved to the wrong room is a different failure from an invented lecturer,
and conflating them hides which half of the pipeline is at fault.

\subsection{Comparison Conventions}

Three are fixed in advance. Source-text spans are excluded, since annotators
disagree about boundaries without disagreeing about the reading. Free-prose
fields on the two open categories are excluded for the same reason: the thing
under test is whether the writer recognised that nothing fits, not whether two
people worded a note alike. Soft weights are excluded by default, a weight
being a scale choice and not a reading, with a weight-inclusive variant
reported as a sensitivity analysis.

\subsection{Keeping the Comparison Fair}

Models from several providers sit behind one interface, and a run is
identified by its provider, model, prompt digest and reasoning effort. Two
provider features that would break symmetry are disabled: native JSON mode,
which the reference configuration does not use, and sampling temperature,
which some of the models do not accept. An instrument that differs between
arms measures the harness. Every raw response is recorded, so a scored run can
be re-scored without further calls.

%% =====================================================================
\section{Results}
\label{sec:ws-results}
%% =====================================================================
\awaitingdata{The corpus has not yet been collected. Subsection headings are
retained to fix the reporting order.}

\subsection{Agreement}
\subsection{Translation Accuracy}
\subsection{Error Categories}
\subsection{Descriptive References}
\subsection{Cross-Provider Comparison}
\subsection{Unsupported and Ambiguous Requests}

%% =====================================================================
\section{Discussion}
%% =====================================================================
\awaitingdata{Depends on Section~\ref{sec:ws-results}.}

%% =====================================================================
\section{Limitations}
%% =====================================================================

Two limitations are structural and hold whatever the corpus shows.

\emph{Accuracy is not usefulness.} A translation measured against a gold label
says nothing about whether a person checking the output ends up better off.
The closest timed comparison of a natural-language interface against a manual
alternative found the interface faster and no more accurate, with correct
output per hour unchanged~\cite{ipeirotis2026nlidb}. That effect is invisible
to the measurement reported here, and a companion study addresses it.

\emph{One vocabulary, one department.} The target vocabulary was frozen early
to keep translation stable, which bounds what staff could express. The open
categories measure that bound instead of hiding it, but a department with
different scheduling practices would need a different vocabulary, and the
accuracy figures would not transfer unchanged.

%% =====================================================================
\section{Conclusion}
%% =====================================================================
\awaitingdata{Depends on Section~\ref{sec:ws-results}.}

\bibliographystyle{ACM-Reference-Format}
\bibliography{references}

\end{document}
